% ============================================================
%  EngineeringMath — shared visual style
%  (palette, headings, theorem family, header/footer, figures,
%   cover and chapter pages, contents)
%
%  Identity: indigo ink, magenta-rose for results, jade for worked
%  examples, amber for exercises — the colours of a domain-colouring
%  plot, on plain white paper.
% ============================================================

% ---------- palette ----------
\definecolor{EMink}{HTML}{1E1F3B}
\definecolor{EMindigo}{HTML}{33307F}
\definecolor{EMindigoL}{HTML}{EDECF8}
\definecolor{EMrose}{HTML}{B3215A}
\definecolor{EMroseL}{HTML}{FAEAF1}
\definecolor{EMjade}{HTML}{0F7B6C}
\definecolor{EMjadeL}{HTML}{E6F4F1}
\definecolor{EMamber}{HTML}{B5700F}
\definecolor{EMamberL}{HTML}{FCF1DE}
\definecolor{EMmuted}{HTML}{6B6B85}
\definecolor{EMviolet}{HTML}{6B3FA0}
\definecolor{EMline}{HTML}{CFCEE2}
\definecolor{EMproofbg}{HTML}{F6F6FB}
\color{EMink}

% ---------- bidi-aware helpers ----------
\providecommand{\lr}[1]{\LR{#1}}
\ifEMrtl
  \newenvironment{EMltr}{\begin{latin}}{\end{latin}}
  \newcommand{\EMnum}[1]{\lr{\persianfont #1}}
  \def\EMstart{east}\def\EMend{west}
  \newcommand{\EMsep}{{\color{EMrose}\rule[0.2em]{0.3em}{0.3em}}}
  \newcommand{\EMheadfont}{\EMfaDisplay}
  \newcommand{\EMmathfont}{}
\else
  \newenvironment{EMltr}{}{}
  \newcommand{\EMnum}[1]{#1}
  \def\EMstart{west}\def\EMend{east}
  \newcommand{\EMsep}{{\color{EMrose}\rule[0.2em]{0.3em}{0.3em}}}
  \newcommand{\EMheadfont}{\sffamily}
\fi
% an inline formula: always left to right, never split
\newcommand{\EMim}[1]{\ifEMrtl\lr{$#1$}\else$#1$\fi}
% Persian words inside a formula run right to left
\newcommand{\EMt}[1]{\mbox{\ifEMrtl\RL{\persianfont #1}\else#1\fi}}
\providecommand{\tightlist}{\setlength{\itemsep}{1pt}\setlength{\parskip}{0pt}}
% symbols missing from Libertinus Math
\usepackage{pifont}
\newcommand{\EMcheck}{\mathord{\text{\textcolor{EMjade}{\ding{51}}}}}
\def\circlearrowright{\mathord{\cwopencirclearrow}}
\def\circlearrowleft{\mathord{\acwopencirclearrow}}
\def\blacksquare{\mathord{\mdlgblksquare}}
\def\square{\mathord{\mdlgwhtsquare}}
\setoperatorfont\symup
% a displayed formula, centred and always left to right.  The filter marks the places where it
% may be broken (before a top-level =, \Rightarrow ..., after \qquad); the macro fills lines
% greedily by their measured width and indents the continuation lines.  A formula that fits
% stays on one line; one that cannot be broken is left to overflow and is reported by the build.
\newtoks\EMdcur
\newbox\EMdb \newbox\EMdvb
\newcount\EMdn \newcount\EMdlines
\newdimen\EMdlim
\newcommand{\EMdline}[1]{\hbox{\ifEMrtl\beginL\fi$\displaystyle#1$\ifEMrtl\endL\fi}}
\def\EMdflush{%
  \setbox\EMdvb\vbox{\offinterlineskip\unvbox\EMdvb
    \ifnum\EMdlines>0 \vskip0.55em\fi
    \hbox{\ifnum\EMdlines>0 \hskip2em\fi\EMdline{\the\EMdcur}}}%
  \advance\EMdlines1 \EMdcur={}\EMdn=0 }
\long\def\EMdseg#1{%
  \ifnum\EMdn>0
    \setbox\EMdb\hbox{$\displaystyle\the\EMdcur #1$}%
    \ifdim\wd\EMdb>\EMdlim \EMdflush \fi
  \fi
  \EMdcur=\expandafter{\the\EMdcur #1}\advance\EMdn1 }
\long\def\EMdproc#1\EMbr#2\EMdend{%
  \EMdseg{#1}%
  \ifblank{#2}{}{\EMdproc#2\EMdend}}
\newcommand{\EMdisp}[1]{\par\nopagebreak[4]\vspace{0.3em}%
  \begingroup\EMdcur={}\EMdn=0 \EMdlines=0 \setbox\EMdvb\vbox{}%
  \EMdlim=\dimexpr\linewidth-2em\relax
  \EMdproc#1\EMbr\EMdend
  \ifnum\EMdlines>0 \EMdflush \else
    \setbox\EMdvb\hbox{\EMdline{\the\EMdcur}}\fi
  \ifdim\wd\EMdvb>\linewidth\setbox\EMdvb\hbox{\resizebox{\linewidth}{!}{\box\EMdvb}}\fi
  \noindent\hfil\box\EMdvb\hfil\par\endgroup\nopagebreak[2]\vspace{0.1em}}
\relpenalty=9999 \binoppenalty=9999
\emergencystretch=3em
\setlength{\parindent}{0pt}
\setlength{\parskip}{0.55em plus 0.15em}
\setlist{itemsep=2pt,topsep=3pt,parsep=0pt}
\setlist[itemize]{label={\color{EMrose}\textbullet}}
\setlist[enumerate,1]{label={\color{EMindigo}\bfseries\arabic*.}}
\setlength{\abovedisplayskip}{0.7em plus 0.2em}
\setlength{\belowdisplayskip}{0.7em plus 0.2em}
\setlength{\abovedisplayshortskip}{0.4em}
\setlength{\belowdisplayshortskip}{0.4em}
\allowdisplaybreaks[1]

% ---------- figures (template/figures/*.tex: TikZ / PGFPlots) ----------
% \EML{English}{فارسی}: a word in a diagram, in the language of the document
\ifEMrtl\newcommand{\EML}[2]{\RL{\persianfont #2}}\else\newcommand{\EML}[2]{#1}\fi
\input{\EMroot/figures/em-diagram-styles}

\newcounter{EMchap}
\newcounter{EMitem}[EMchap]
\newcounter{EMfig}[EMchap]
\renewcommand{\thesection}{\arabic{EMchap}.\arabic{section}}
% unique PDF anchors across chapters; chapters are bookmark level 1, sections level 2
\renewcommand{\theHsection}{\arabic{EMchap}.\arabic{section}}
\renewcommand{\theHEMitem}{\arabic{EMchap}.\arabic{EMitem}}
\renewcommand{\theHEMfig}{\arabic{EMchap}.\arabic{EMfig}}
\makeatletter\renewcommand*{\toclevel@section}{2}\makeatother
% Persian digits for a bookmark title (expandable)
\newcommand{\EMbmnum}[1]{\ifEMrtl\ifcase#1 ۰\or ۱\or ۲\or ۳\or ۴\or ۵\or ۶\or ۷\or ۸\or ۹\or ۱۰\or ۱۱\or ۱۲\or ۱۳\or ۱۴\else #1\fi\else#1\fi}
\newcommand{\EMitemnum}{\EMnum{\arabic{EMchap}.\arabic{EMitem}}}
\newcommand{\EMfignum}{\EMnum{\arabic{EMchap}.\arabic{EMfig}}}

\newsavebox{\EMfigbox}
\newlength{\EMfigh}
% the drawing, scaled down only when it would exceed the text width
\newcommand{\EMfigmeasure}[1]{%
  \sbox{\EMfigbox}{\begin{EMltr}\input{\EMroot/figures/#1}\end{EMltr}}%
  \ifdim\wd\EMfigbox>\linewidth
    \sbox{\EMfigbox}{\resizebox{\linewidth}{!}{\usebox{\EMfigbox}}}\fi
  \setlength{\EMfigh}{\dimexpr\ht\EMfigbox+\dp\EMfigbox\relax}}
% reserves the room of a figure (and the lines before it) so a heading is never left alone
\newcommand{\EMkeepfig}[2]{\par\begingroup\EMfigmeasure{#1}%
  \Needspace*{\dimexpr\EMfigh+\numexpr#2+4\relax\baselineskip\relax}\endgroup}
% #1 = figure file, #2 = caption; drawing and caption never split
\newcommand{\EMfigure}[2]{%
  \par\EMfigmeasure{#1}\Needspace*{\dimexpr\EMfigh+4\baselineskip\relax}%
  \refstepcounter{EMfig}\label{fig:#1}\vspace{0.5em}%
  \begin{minipage}{\linewidth}\centering\usebox{\EMfigbox}\par\vspace{0.5em}%
  {\small\color{EMmuted}{\color{EMindigo}\EMheadfont\bfseries\EMlangFigure~\EMfignum}\hspace{0.7em}#2\par}%
  \end{minipage}\par\vspace{0.8em}}

% ---------- headings ----------
\newcommand{\EMmodkey}{x}
\newcommand{\EMsectag}{\ifEMrtl\lr{\fi\tikz[baseline=(n.base)]\node[fill=EMindigo,text=white,
  rounded corners=1.8pt,inner xsep=4.5pt,inner ysep=3pt,font=\EMheadfont\bfseries\small](n){\EMnum{\thesection}};\ifEMrtl}\fi}
\newcommand{\EMsectionreal}[2]{%
  \par\Needspace{9\baselineskip}\vspace{1.3em}%
  \refstepcounter{section}\label{EMsec:\EMmodkey:\arabic{section}}%
  \addcontentsline{toc}{section}{\texorpdfstring{#1}{#2}}%
  \noindent\EMsectag\hspace{0.7em}{\EMheadfont\bfseries\Large\color{EMindigo}\hyphenpenalty=10000 \exhyphenpenalty=10000 #1}\par\nobreak
  \vspace{0.25em}{\color{EMline}\hrule height 0.6pt}\vspace{0.6em}\nobreak}
\newcommand{\EMsubsectionreal}[2]{%
  \par\Needspace{6\baselineskip}\vspace{0.9em}%
  \noindent{\color{EMrose}\rule[0.1em]{0.32em}{0.78em}}\hspace{0.55em}%
  {\EMheadfont\bfseries\large\color{EMink}\hyphenpenalty=10000 \exhyphenpenalty=10000 #1}\par\nobreak\vspace{0.25em}\nobreak}

% a heading followed by a box (or a unit, see below) is held back until that block knows how
% much room it needs; the heading and the block then move together
\newif\ifEMhold
\def\EMpending{}\newcount\EMpendn
\newcommand{\EMhold}{\global\EMholdtrue}
\newcommand{\EMqueue}[1]{\gappto\EMpending{#1}\global\advance\EMpendn by 1\relax}
\newcommand{\EMsection}[2]{\ifEMhold\EMqueue{\EMsectionreal{#1}{#2}}\else\EMsectionreal{#1}{#2}\fi}
\newcommand{\EMsubsection}[2]{\ifEMhold\EMqueue{\EMsubsectionreal{#1}{#2}}\else\EMsubsectionreal{#1}{#2}\fi}
\newcommand{\EMflush}{\par\ifnum\EMpendn>0
  \begingroup\global\EMholdfalse\EMpending\endgroup
  \global\let\EMpending\empty\global\EMpendn=0 \fi\global\EMholdfalse}

% ---------- keeping a logical block on one page ----------
% A block (a box, or a "unit": a paragraph with its equations, a figure with its lead-in)
% is typeset once into a box to measure it.  If it fits on a page it is kept whole: when it
% does not fit in what is left of the current page it moves to the next page.  A block longer
% than a page stays breakable and starts with room for its head and a few lines.
\newdimen\EMpendht
\newdimen\EMboxht
\newdimen\EMkeepmax
\newcount\EMsvitem \newcount\EMsvfig
\newcommand{\EMsavecnt}{\EMsvitem=\value{EMitem}\relax\EMsvfig=\value{EMfig}\relax}
\newcommand{\EMrestorecnt}{\setcounter{EMitem}{\EMsvitem}\setcounter{EMfig}{\EMsvfig}}
% #1 = width of the measured text, #2 = what is measured, #3 = extra height (padding)
\newcommand{\EMmeasureraw}[3]{\EMsavecnt
  \setbox0\vbox{\hsize#1\linewidth\hsize\columnwidth\hsize\parindent0pt\relax
    \setlength{\parskip}{0.45em}#2\par}%
  \EMrestorecnt
  \EMboxht=\dimexpr\ht0+\dp0+#3\relax}
\newif\ifEMnoplace
\newcommand{\EMplace}{% decides the room, then places the held headings
  \ifEMnoplace\def\EMboxopt{unbreakable}\else
  \EMkeepmax=\dimexpr\textheight-1.5\baselineskip\relax
  \EMpendht=\dimexpr\EMpendn\baselineskip*3\relax
  \ifdim\EMboxht<\EMkeepmax
    \Needspace*{\dimexpr\EMboxht+\EMpendht+0.6\baselineskip\relax}\def\EMboxopt{unbreakable}%
  \else\Needspace*{\dimexpr7\baselineskip+\EMpendht\relax}\def\EMboxopt{}\fi
  \EMflush\fi}
% a join: a statement and the proof that follows it are one block when together they fit on a page
\NewEnviron{EMjoin}{%
  \EMnoplacetrue\EMmeasureraw{\linewidth}{\BODY}{0pt}\EMnoplacefalse
  \ifdim\EMboxht<\dimexpr\textheight-1.5\baselineskip\relax
    \EMplace\EMnoplacetrue\BODY\EMnoplacefalse
  \else\EMplace\BODY\fi\par}
% a unit: ordinary text that belongs together (generated by the filter)
\NewEnviron{EMunit}{%
  \EMmeasureraw{\linewidth}{\BODY}{0pt}%
  \EMplace
  \BODY\par}

% ---------- the theorem family ----------
\ifEMrtl
  \tcbset{EMbar/.style={borderline east={2.8pt}{0pt}{#1}}}
\else
  \tcbset{EMbar/.style={borderline west={2.8pt}{0pt}{#1}}}
\fi
% one box style: a thin hairline all round in a soft tone of the colour, a stronger bar on the
% leading edge, a very light tint, a hairline under the heading
\tcbset{EMbase/.style={enhanced,breakable,sharp corners,boxrule=0.5pt,
  left=9pt,right=9pt,top=5pt,bottom=6pt,before skip=0.9em,after skip=0.9em,
  before upper={\parindent0pt\setlength{\parskip}{0.45em}},
  overlay unbroken and first={},
  toprule at break=0pt,bottomrule at break=0pt}}
\newcommand{\EMrule}[1]{\par\nobreak\vspace{0.2em}{\color{#1!35}\hrule height 0.4pt}\vspace{0.4em}\nobreak}
\newcommand{\EMboxhead}[3]{%   colour, label word, title (may be empty)
  \refstepcounter{EMitem}\label{EMit:\EMmodkey:\arabic{EMitem}}%
  {\EMheadfont\bfseries\color{#1}#2~\EMitemnum}%
  \if\relax\detokenize{#3}\relax\else{\ }\EMsep{\ }{\EMheadfont\color{#1}#3}\fi\EMrule{#1}}
\newcommand{\EMboxheadplain}[2]{{\EMheadfont\bfseries\color{#1}#2}\EMrule{#1}}
\newcommand{\EMboxheadnn}[3]{{\EMheadfont\bfseries\color{#1}#2}%
  \if\relax\detokenize{#3}\relax\else{\ }\EMsep{\ }{\EMheadfont\color{#1}#3}\fi\EMrule{#1}}
% a numbered or plain box.  #1 env, #2 colour, #3 tint, #4 label word, #5 numbered (1/0)
\newcommand{\EMmakebox}[5]{%
  \NewEnviron{#1}[1]{%
    \EMmeasureraw{\dimexpr\linewidth-20pt\relax}{%
      \if#51\EMboxhead{#2}{#4}{##1}\else\EMboxheadnn{#2}{#4}{##1}\fi\BODY}{16pt}%
    \EMplace
    \begin{tcolorbox}[EMbase,\EMboxopt,EMbar=#2,colback=#3,colframe=#2!40!white]%
    \if#51\EMboxhead{#2}{#4}{##1}\else\EMboxheadnn{#2}{#4}{##1}\fi
    \BODY\end{tcolorbox}}}
\EMmakebox{EMdefinition}{EMindigo}{EMindigoL}{\EMlangDefinition}{1}
\EMmakebox{EMtheorem}{EMrose}{EMroseL}{\EMlangTheorem}{1}
\EMmakebox{EMlemma}{EMrose}{EMroseL}{\EMlangLemma}{1}
\EMmakebox{EMproposition}{EMrose}{EMroseL}{\EMlangProposition}{1}
\EMmakebox{EMcorollary}{EMrose}{EMroseL}{\EMlangCorollary}{1}
\EMmakebox{EMexample}{EMjade}{EMjadeL}{\EMlangExample}{1}
\EMmakebox{EMexercise}{EMamber}{EMamberL}{\EMlangExercise}{1}
\EMmakebox{EMremark}{EMmuted}{white}{\EMlangRemark}{0}
% the key results: the same box, a stronger frame
\NewEnviron{EMimportant}[1]{%
  \EMmeasureraw{\dimexpr\linewidth-20pt\relax}{\EMboxhead{EMrose}{\EMlangImportant}{#1}\BODY}{18pt}%
  \EMplace
  \begin{tcolorbox}[EMbase,\EMboxopt,EMbar=EMrose,colback=EMroseL,colframe=EMrose,boxrule=0.9pt,
    top=6pt,bottom=7pt]%
  \EMboxhead{EMrose}{\EMlangImportant}{#1}\BODY\end{tcolorbox}}
% a proof: a quiet box with a clear heading and a closing mark at the end of the last line
\newcommand{\EMqed}{\par\nobreak\vspace{-0.15em}\nopagebreak{\hfill\color{EMindigo}\ensuremath{\blacksquare}}\par}
\NewEnviron{EMproof}[1]{%
  \EMmeasureraw{\dimexpr\linewidth-20pt\relax}{%
    \EMboxheadplain{EMindigo}{\EMlangProof.}\BODY\EMqed}{16pt}%
  \EMplace
  \begin{tcolorbox}[EMbase,\EMboxopt,EMbar=EMindigo,colback=EMproofbg,colframe=EMindigo!30!white]%
  \EMboxheadplain{EMindigo}{\EMlangProof.}\BODY\EMqed\end{tcolorbox}}
% a solution inside an example or an exercise: a run-in heading
\newenvironment{EMsolution}{\par\noindent{\EMheadfont\bfseries\color{EMjade}\EMlangSolution.}\ }{\par}

% ---------- header / footer ----------
\def\EMchaptitle{}\def\EMchapother{}
\newcommand{\EMsetmodule}[3]{\gdef\EMmodkey{#1}\gdef\EMchaptitle{#2}\gdef\EMchapother{#3}%
  \ifnum0#1>0 \setcounter{EMchap}{#1}\fi\setcounter{section}{0}}
\newcommand{\EMheadchapter}{\ifnum\value{EMchap}>0 {\color{EMrose}\EMlangChapter~\EMnum{\arabic{EMchap}}}\ \EMsep\ \fi\EMchaptitle}
\newcommand{\EMheadcourse}{\EMlangCourse}
\pagestyle{fancy}
\fancyhf{}
\fancyhead[L]{\small\color{EMmuted}\ifEMrtl\EMheadchapter\else\EMheadcourse\fi}
\fancyhead[R]{\small\color{EMmuted}\ifEMrtl\EMheadcourse\else\EMheadchapter\fi}
\renewcommand{\headrulewidth}{0pt}
\fancyfoot[C]{\lr{\tikz[baseline=(p.base)]\node[circle,draw=EMindigo,line width=0.7pt,inner sep=2.2pt,
  minimum size=1.9em,font=\small\EMheadfont,text=EMindigo](p){\EMnum{\thepage}};}}
\fancypagestyle{plain}{\fancyhf{}\renewcommand{\headrulewidth}{0pt}%
  \fancyfoot[C]{\lr{\tikz[baseline=(p.base)]\node[circle,draw=EMindigo,line width=0.7pt,inner sep=2.2pt,
  minimum size=1.9em,font=\small\EMheadfont,text=EMindigo](p){\EMnum{\thepage}};}}}

% ---------- contents ----------
\newcommand{\EMtocleader}{\leaders\hbox{\kern1.5pt\textcolor{EMline}{.}\kern1.5pt}\hfill}
\newenvironment{EMminitoc}{%
  \par\vspace{0.6em}{\EMheadfont\bfseries\large\color{EMindigo}\EMlangChapterContents}\par\vspace{0.3em}%
  {\color{EMline}\hrule height 0.6pt}\vspace{0.5em}\par}{\par}
\newcommand{\EMtocitem}[2]{% section number in the chapter, title
  \par\noindent\hyperref[EMsec:\EMmodkey:#1]{\hbox to\linewidth{\EMsectagtoc{#1}\hspace{0.6em}#2\EMtocleader\hspace{0.4em}\pageref{EMsec:\EMmodkey:#1}}}\par}
\newcommand{\EMsectagtoc}[1]{\ifEMrtl\lr{\fi\tikz[baseline=(n.base)]\node[fill=EMindigo,text=white,
  rounded corners=1.5pt,inner xsep=3.5pt,inner ysep=2pt,font=\EMheadfont\bfseries\footnotesize](n){\EMnum{\arabic{EMchap}.#1}};\ifEMrtl}\fi}
% the book: parts, chapters, sections
\newenvironment{EMbooktoc}{%
  \par{\EMheadfont\bfseries\fontsize{26}{30}\selectfont\color{EMindigo}\EMlangContents}\par\vspace{0.4em}%
  {\color{EMline}\hrule height 0.8pt}\vspace{0.8em}\par}{\par}
\newcommand{\EMtocpart}[2]{\par\Needspace{6\baselineskip}\vspace{1.1em}\noindent
  {\EMheadfont\bfseries\large\color{EMrose}\EMlangPart~#1}\ \ \EMsep\ \ {\EMheadfont\bfseries\large\color{EMink}#2}\par\vspace{0.1em}}
\newcommand{\EMtocchapter}[3]{% key, number, title
  \par\Needspace{4\baselineskip}\vspace{0.5em}\noindent\hyperref[EMch:#1]{\hbox to\linewidth{%
    \EMheadfont\bfseries\color{EMindigo}\EMnum{#2}\hspace{0.7em}#3\EMtocleader\hspace{0.4em}\pageref{EMch:#1}}}\par}
\newcommand{\EMtocitemk}[4]{% key, k, number, title
  \par\noindent\hspace{1.4em}\hyperref[EMsec:#1:#2]{\hbox to\dimexpr\linewidth-1.4em\relax{\small
    \EMnum{#3}\hspace{0.6em}#4\EMtocleader\hspace{0.4em}\pageref{EMsec:#1:#2}}}\par}

% ---------- cover motif: a Joukowski map of a polar grid ----------
% the circles |z|=r map to ellipses, the rays arg z=t to hyperbola branches
\newcommand{\EMjoukowski}[3]{% scale, indigo strength, rose strength
  \begin{scope}[x=#1,y=#1,line cap=round]
    \foreach \r in {1.06,1.14,1.26,1.42,1.65,2,2.6,3.6}{%
      \draw[EMindigo!#2,line width=0.5pt] plot[domain=0:360,samples=90,smooth]
        ({(\r+1/\r)*cos(\x)},{(\r-1/\r)*sin(\x)});}
    \foreach \t in {10,25,...,175}{%
      \draw[EMrose!#3,line width=0.45pt] plot[domain=1:4,samples=40,smooth]
        ({(\x+1/\x)*cos(\t)},{(\x-1/\x)*sin(\t)});
      \draw[EMrose!#3,line width=0.45pt] plot[domain=1:4,samples=40,smooth]
        ({(\x+1/\x)*cos(\t)},{-(\x-1/\x)*sin(\t)});}
    \draw[EMink!80,line width=0.8pt] (-2,0)--(2,0);
  \end{scope}}

% ---------- motif styles (cover and chapter openers) ----------
\tikzset{
  max/.style={-{Stealth[length=5pt,width=4pt]},line width=0.6pt,EMink!70},
  mfaint/.style={line width=0.35pt,EMline},
  mline/.style={line width=1.1pt,EMindigo},
  mrose/.style={line width=1.1pt,EMrose},
  mjade/.style={line width=1.1pt,EMjade},
  mdash/.style={dashed,line width=0.6pt,EMmuted},
  mvec/.style={-{Stealth[length=6pt,width=5pt]},line width=1.2pt,EMrose},
  mvecl/.style={-{Stealth[length=6pt,width=5pt]},line width=1.2pt,EMindigo!75},
  mmid/.style={decoration={markings,mark=at position 0.5 with {\arrow{Stealth[length=6pt]}}},postaction=decorate},
  mdot/.style={circle,inner sep=0pt,minimum size=5pt},
  mdotr/.style={circle,fill=EMrose,inner sep=0pt,minimum size=5.5pt},
  mdotl/.style={circle,fill=EMindigo,inner sep=0pt,minimum size=5.5pt},
}
\newcommand{\EMcrossm}[1]{\draw[EMrose,line width=1.4pt] #1 ++(-3.5pt,-3.5pt)--++(7pt,7pt) #1 ++(-3.5pt,3.5pt)--++(7pt,-7pt);}
% the motif of chapter #1 (a 15 cm x 9 cm drawing)
\newcommand{\EMmotif}[1]{\input{\EMroot/motifs/m\ifnum#1<10 0\fi#1}}
% Joukowski map on a dark ground: colours given as full colour specs
\newcommand{\EMjoukowskiD}[3]{% scale, circle colour, ray colour
  \begin{scope}[x=#1,y=#1,line cap=round]
    \foreach \r in {1.06,1.14,1.26,1.42,1.65,2,2.6,3.6}{%
      \draw[#2,line width=0.5pt] plot[domain=0:360,samples=90,smooth]
        ({(\r+1/\r)*cos(\x)},{(\r-1/\r)*sin(\x)});}
    \foreach \t in {10,25,...,175}{%
      \draw[#3,line width=0.45pt] plot[domain=1:4,samples=40,smooth]
        ({(\x+1/\x)*cos(\t)},{(\x-1/\x)*sin(\t)});
      \draw[#3,line width=0.45pt] plot[domain=1:4,samples=40,smooth]
        ({(\x+1/\x)*cos(\t)},{-(\x-1/\x)*sin(\t)});}
    \draw[white!80,line width=0.8pt] (-2,0)--(2,0);
  \end{scope}}

% ---------- cover ----------
% Structure after the ProbAndStats cover (university band, logo, title block, ornament, card,
% credits), with the Engineering Mathematics art panel at the foot.  The page is laid out in
% the direction of its language: in the Persian edition every Persian string is an RTL run
% (\RL), the leading edge is the right-hand one, and the credits read right to left.
\ifEMrtl
  \newcommand{\EMprim}[1]{\RL{\persianfont #1}}          % text in the language of the edition
  \newcommand{\EMoth}[1]{\lr{#1}}                        % text in the other language
  \newcommand{\EMprimD}[1]{\RL{\EMfaDisplay #1}}          % display weight
  \def\EMcovStart{east}\def\EMcovEnd{west}
\else
  \newcommand{\EMprim}[1]{#1}
  \newcommand{\EMoth}[1]{\rl{#1}}
  \newcommand{\EMprimD}[1]{#1}
  \def\EMcovStart{west}\def\EMcovEnd{east}
\fi
% the art panel: a Joukowski map, a contour round three poles, a decaying oscillation and the
% harmonics of a square wave, on the dark ground at the foot of the page
\newcommand{\EMcoverart}{%
  \begin{scope}[shift={($(current page.south west)+(10.5,4.9)$)}]
    \EMjoukowskiD{0.82cm}{EMindigo!55!white!70}{EMrose!85!white!80}\end{scope}
  \begin{scope}[shift={($(current page.south west)+(3.3,5.0)$)},scale=0.62]
    \draw[white!75,line width=1.1pt,mmid,smooth cycle,tension=0.85] plot coordinates{(-2.6,-0.2)(-1.6,1.8)(0.8,2.3)(2.6,0.9)(2.2,-1.6)(-0.3,-2.2)};
    \foreach \p in {(-0.9,0.5),(0.9,-0.4),(0.5,1.2)}{\draw[EMrose!70!white,line width=0.9pt] \p circle (0.3);\fill[EMrose!60!white] \p circle (1.6pt);}
  \end{scope}
  \begin{scope}[shift={($(current page.south west)+(15.3,4.5)$)},scale=0.62]
    \draw[white!50,line width=0.5pt,-{Stealth[length=5pt]}] (0,0)--(5.4,0);
    \draw[white!35,dashed,line width=0.5pt,domain=0:5.2,samples=50] plot (\x,{1.9*exp(-0.5*\x)});
    \draw[EMjade!60!white,line width=1.4pt,domain=0:5.2,samples=250] plot (\x,{1.9*exp(-0.5*\x)*cos(4.8*\x r)});
  \end{scope}
  \begin{scope}[shift={($(current page.south west)+(0,1.25)$)}]
    \draw[white!30,line width=0.4pt] (0,0)--(21,0);
    \draw[EMjade!70!white,line width=0.6pt,domain=0:21,samples=200] plot (\x,{0.5*sin(0.7*1.5*\x r)});
    \draw[EMamber!70!white,line width=0.6pt,domain=0:21,samples=300] plot (\x,{0.5/3*sin(3*0.7*1.5*\x r)});
    \draw[EMrose!50!white,line width=1.6pt,domain=0:21,samples=500] plot (\x,{0.62*(sin(0.7*1.5*\x r)+sin(3*0.7*1.5*\x r)/3+sin(5*0.7*1.5*\x r)/5+sin(7*0.7*1.5*\x r)/7)});
  \end{scope}}

% #1 course title, #2 subtitle (the three parts), #3 title in the other language,
% #4 subtitle in the other language
\newcommand{\EMcoverpage}[4]{%
  \thispagestyle{empty}%
  \begin{tikzpicture}[remember picture,overlay]
    \coordinate (sw) at (current page.south west);
    % the art panel at the foot
    \fill[EMink] (sw) rectangle ([yshift=8.4cm]current page.south east);
    \begin{scope}\clip (sw) rectangle ([yshift=8.4cm]current page.south east);\EMcoverart\end{scope}
    \fill[EMrose] ([yshift=8.4cm]sw) rectangle ([yshift=8.55cm]current page.south east);
    % university band: Persian name on the right, English name on the left
    \fill[EMink] (current page.north west) rectangle ([yshift=-1.25cm]current page.north east);
    \fill[EMrose] ([yshift=-1.25cm]current page.north west) rectangle ([yshift=-1.4cm]current page.north east);
    \node[text=white,font=\small\bfseries,anchor=east] at ([xshift=-1.6cm,yshift=-0.62cm]current page.north east)
      {\RL{\EMfaDisplay دانشگاه فردوسی مشهد}};
    \node[text=white,font=\small\scshape,anchor=west] at ([xshift=1.6cm,yshift=-0.62cm]current page.north west)
      {Ferdowsi University of Mashhad};
    % logo (aspect ratio kept)
    \node[anchor=north] (logo) at ([yshift=-2.35cm]current page.north)
      {\includegraphics[width=3.3cm]{\EMroot/ferdowsi-logo.png}};
    % course title in the language of the edition, the other language beneath it
    \node[anchor=north,text=EMindigo,font=\ifEMrtl\fontsize{46}{56}\selectfont\else\EMheadfont\bfseries\fontsize{40}{48}\selectfont\fi] (title)
      at ([yshift=-0.85cm]logo.south) {\EMprimD{#1}};
    \node[anchor=north,text=EMmuted,font=\fontsize{19}{24}\selectfont] (other)
      at ([yshift=-0.3cm]title.south) {\EMoth{#3}};
    % ornament
    \draw[EMrose,line width=0.9pt] ([xshift=-4.2cm,yshift=-0.6cm]other.south) -- ([xshift=-0.35cm,yshift=-0.6cm]other.south);
    \draw[EMrose,line width=0.9pt] ([xshift=0.35cm,yshift=-0.6cm]other.south) -- ([xshift=4.2cm,yshift=-0.6cm]other.south);
    \fill[EMrose] ([yshift=-0.6cm]other.south) +(0,0.16) -- +(0.16,0) -- +(0,-0.16) -- +(-0.16,0) -- cycle;
    % card: what the notes cover
    \node[anchor=north,fill=EMindigo,rounded corners=4mm,inner xsep=9mm,inner ysep=5mm,
          text width=14.2cm,align=center,text=white] (card) at ([yshift=-1.15cm]other.south)
      {{\color{EMrose!55!white}\large\bfseries\EMprim{\EMlangCompleteNotes}}\\[2.5mm]
       {\ifEMrtl\fontsize{15}{22}\else\fontsize{12.5}{20}\fi\selectfont\bfseries\hyphenpenalty=10000 \EMprim{#2}}\\[2.5mm]
       {\small\color{white!78!EMindigo}\EMoth{#4}}};
    % instructor and compiler, in the direction of the edition
    \node[anchor=north,text=EMink,align=center] at ([yshift=-0.85cm]card.south)
      {{\large\bfseries\EMprim{\EMlangInstructorLabel:\ \EMlangInstructor}}\\[1.4mm]
       {\normalsize\color{EMindigo}\EMprim{\EMlangCompilerLabel:\ \EMlangCompiler}}\\[2.2mm]
       {\small\color{EMmuted}\EMoth{\EMlangInstructorLineOther}}};
  \end{tikzpicture}\clearpage}

% ---------- chapter opener (also the first page of a single-chapter PDF) ----------
% #1 key, #2 number, #3 title, #4 title in the other language.  The chapter contents follow
% on the same page (\EMchapterend closes it); the text starts on the next page.
\newcommand{\EMchapterpage}[4]{%
  \clearpage\pdfbookmark[1]{\EMlangChapter~\EMbmnum{#2}: #3}{EMch@#1}\phantomsection\label{EMch:#1}%
  \thispagestyle{plain}%
  \begin{tikzpicture}[remember picture,overlay]
    \fill[EMindigo] ([xshift=\ifEMrtl-1.1cm\else0cm\fi]current page.north \ifEMrtl east\else west\fi) rectangle
      ([xshift=\ifEMrtl0cm\else1.1cm\fi]current page.south \ifEMrtl east\else west\fi);
  \end{tikzpicture}%
  \vspace*{-0.6cm}%
  \par\noindent\hbox to\linewidth{{\EMheadfont\bfseries\large\color{EMmuted}\EMlangChapter}\hfill}\par\vspace{-0.5em}
  \noindent\hbox to\linewidth{{\EMheadfont\bfseries\fontsize{96}{100}\selectfont\color{EMrose!85}\EMnum{#2}}\hfill}\par\vspace{0.1em}
  {\color{EMrose}\rule{3.2cm}{2pt}}\par\vspace{0.7em}
  {\EMheadfont\bfseries\fontsize{25}{31}\selectfont\color{EMindigo}\hyphenpenalty=10000 \exhyphenpenalty=10000 \raggedright#3\par}
  \vspace{0.3em}{\large\color{EMmuted}\EMother{#4}\par}
  \vspace{0.4em}
  \begin{center}
    \begin{tikzpicture}[scale=0.78]
      \clip (0,0) rectangle (15,9);
      \EMmotif{#1}
    \end{tikzpicture}
  \end{center}}
\newcommand{\EMchapterend}{\clearpage}

% ---------- part page (the book) ----------
\newcommand{\EMpartpage}[5]{% key, numeral, title, other title, chapter range
  \clearpage\pdfbookmark[0]{\EMlangPart~#2: #3}{EMpart@#1}\phantomsection\label{EMpart:#1}%
  \thispagestyle{empty}%
  \begin{tikzpicture}[remember picture,overlay]
    \begin{scope}[shift={($(current page.center)+(0,3.2)$)}]\EMjoukowski{1.02cm}{50}{60}\end{scope}
    \fill[EMindigo] ([yshift=8.2cm]current page.south west) rectangle ([yshift=2.4cm]current page.south east);
    \node[anchor=south,text=EMline,align=center,text width=16cm,font=\EMheadfont\large]
      at ($(current page.south)+(0,6.55)$) {\EMlangPart\ #2};
    \node[anchor=south,text=white,align=center,text width=16cm,font=\EMheadfont\bfseries\fontsize{30}{36}\selectfont]
      at ($(current page.south)+(0,4.7)$) {#3};
    \node[anchor=south,text=white!75!EMindigo,align=center,text width=16cm,font=\large]
      at ($(current page.south)+(0,3.55)$) {\EMother{#4}};
    \node[anchor=south,text=EMline,align=center,text width=16cm,font=\small]
      at ($(current page.south)+(0,2.75)$) {#5};
  \end{tikzpicture}\clearpage}

% ---------- about these notes ----------
\newcommand{\EMcolophon}{%
  \clearpage\thispagestyle{plain}%
  {\EMheadfont\bfseries\fontsize{22}{26}\selectfont\color{EMindigo}\EMlangColophonTitle\par}
  \vspace{0.3em}{\color{EMline}\hrule height 0.8pt}\vspace{0.8em}
  \EMcolophontext
  \vspace{0.6em}
  \begin{EMdefinition}{\EMlegendDefinition}\EMlegendDefinitionText\end{EMdefinition}
  \begin{EMtheorem}{\EMlegendTheorem}\EMlegendTheoremText\end{EMtheorem}
  \begin{EMimportant}{\EMlegendImportant}\EMlegendImportantText\end{EMimportant}
  \begin{EMexample}{\EMlegendExample}\EMlegendExampleText
    \begin{EMsolution}\EMlegendSolutionText\end{EMsolution}\end{EMexample}
  \begin{EMexercise}{\EMlegendExercise}\EMlegendExerciseText\end{EMexercise}
  \begin{EMremark}{\EMlegendRemark}\EMlegendRemarkText\end{EMremark}
  \setcounter{EMitem}{0}\clearpage}
